% GENERATED by experiments/analyze_klsft.py -- do not edit by hand.
\newcommand{\KLBetas}{0.5, 1.0}
\newcommand{\KLNSeeds}{5}
\newcommand{\KLBetaHalfTransferDelta}{$+0.061$}
\newcommand{\KLBetaHalfRepDelta}{$-0.035$}
\newcommand{\KLBetaOneTransferDelta}{$+0.051$}
\newcommand{\KLBetaOneRepDelta}{$-0.053$}
\newcommand{\KLVerdictSentence}{In other words, a base-anchored penalty recovers a meaningful part of the transfer that vanilla SFT gives up --- so a substantial share of Act~I's transfer loss is a property of the \emph{unregularized} recipe, not of fine-tuning as such.}
\newcommand{\KLTakeaway}{A one-line change to the recipe --- adding $\beta\,\mathrm{KL}(\pi_\theta\|\pi_{\text{base}})$ --- recovers much of the transfer that plain SFT sacrifices, at a modest represented cost. The Act~I specialization is therefore partly a property of the recipe: it is mitigable \emph{within} the SFT family, without the composition of Act~II. \textbf{Two limits travel with that, and both point the same way.} This is a retrospective estimate on four general checkpoints with no interval attached; the \emph{preregistered} test in \Cref{sec:adaptation} finds the represented-source cost of the same trade fails its non-inferiority margin (RQ2 \textbf{not supported}). And on the two checkpoints that specialize hardest, KL-SFT still leaves held-out transfer \emph{below} the unmodified base --- mitigation, not restoration. Treat $\beta$ as a tradeoff dial, not a default.}
